\begin{textsample}{POS Dim 1 – persona\_gemini – Score 13.00 – t704\_gemini.txt}  \label{ex:f1_pos_046}
On August 19, we celebrate World Photography Day, and it brings to \textbf{mind} Robert Capa 's advice : if your \textbf{pictures} aren't \textbf{good} enough, you 're not close enough.
As Greenpeace International 's Multimedia \textbf{Editor} in Hong Kong, I see how photography has a profound nature that is simultaneously personal and public.
Photography 's origins are in the 19th \textbf{century}, \textbf{bearing} \textbf{witness} through landscapes and portraits.
For over 200 \textbf{years}, \textbf{images} have served as evidence in court and have inspired action.
As a powerful ally to environmental \textbf{movements}, photography communicates universally.
This \textbf{post} features 12 iconic Greenpeace \textbf{images}.
They include a turtle entangled in plastic, a bloodied \textbf{whale} which helped prompt whaling bans, oil-soaked pelicans, and \textbf{scenes} of deforestation in Indonesia.
Also featured are \textbf{images} of protesters at Standing Rock and a \textbf{woman} victorious in her fight against a coal company.
We urge you to join the \textbf{movement}.

% matched lemmas: bear, century, editor, good, image, mind, movement, picture, post, scene, whale, witness, woman, year
\end{textsample}
